\chapter{The database: status, migration, backup}\label{ch:database}

Every Test Run is a record in one SQLite database, at \texttt{database.path}. ATE writes it as the
run happens and never edits or deletes a run afterwards: v1 has no delete and no purge, so the
database is an archive, not a scratch pad.

\section{The two commands}

\begin{lstlisting}[language=ATEShell]
ate db status
ate db migrate
ate db migrate --yes
ate --config /tmp/station/config/config.json db status
ate db status --config /tmp/station/config/config.json
ate db migrate /tmp/station/data/ate.db --yes
\end{lstlisting}

\texttt{ate db status} opens the database read-only and prints its schema version, exiting 0. A
missing, empty or older database exits 1 and prints the refusal. \texttt{ate db migrate} creates or
upgrades it. Without \texttt{--yes} (or \texttt{-y}) it asks first:

\begin{lstlisting}[language=ATEShell]
Create or upgrade database at /var/lib/ate/ate.db? [y/N]
\end{lstlisting}

Answering anything other than \texttt{y} prints \texttt{aborted} and exits 1. With \texttt{--yes}
there is no prompt, which is what scripts and \texttt{ate setup} use. Neither command needs a
config: give it a path, or let it use the configured database from the platform config.
\texttt{--config} belongs to \texttt{ate}, and it may sit \emph{before} the subcommand
(\texttt{ate --config PATH db status}) or \emph{after} it (\texttt{ate db status --config PATH}):
either reads that config. Written after the subcommand it is no longer read as a database path,
which is the trap this build closes. A path given after \texttt{db migrate} still wins over the
configured one. A path that is not a database is refused in one line with exit 1, never with a
traceback:

\begin{lstlisting}[language=ATEShell]
the file at /tmp/station/data/notes.db is not a SQLite database
cannot open the database at /var/lib/ate/ate.db: unable to open database file
\end{lstlisting}

The first is a file that is not a database at all; the second is a real database this process cannot
open. They are different faults, and the fix is different: replace the file, or fix the permissions
and the ownership of the directory it sits in.

Migrations are forward only. A database built by a \emph{newer} ATE is refused rather than
downgraded; a missing, empty or older one is exactly what migrate is for. The version is SQLite's
\texttt{PRAGMA user\_version}, and the current one is 5. A database that stores an older version but
already holds every table this ATE declares is also refused, naming the version it stores:

\begin{lstlisting}[language=ATEShell]
the database at /var/lib/ate/ate.db stores schema version 4 but is missing nothing to migrate:
no step from version 4 to 5 ran, so it is not stamped as version 5
\end{lstlisting}

Nothing would be created, so stamping it as current would be a claim with no migration behind it.

\section{What the app itself does at launch}

The application never migrates. It checks the version and stops if it is wrong:

\begin{lstlisting}[language=ATEShell]
run `ate db migrate` to create or upgrade the database at /var/lib/ate/ate.db
this ATE is too old for the database at /var/lib/ate/ate.db; update ATE
\end{lstlisting}

The first line names the database, and when the Station is reading a config that is not the
platform default the hint carries \texttt{--config <path>} so you migrate the right one.

At launch ATE also closes any run that has no end time --- a run left open by a crash or a power
cut --- recording it as \texttt{INTERRUPTED} with result \texttt{FAIL}, and logs one line saying
how many it closed: \texttt{interrupted run(s) at launch}. That is a record of what was lost, not a
bug to fix.

\section{Backups}

A migration takes its own copy before it changes anything. \texttt{ate db migrate} and
\texttt{ate setup} write a copy beside the database, named
\texttt{ate.db.backup-<YYYYMMDD-HHMMSS>} in UTC, and print where it went:

\begin{lstlisting}[language=ATEShell]
backed up /var/lib/ate/ate.db to /var/lib/ate/ate.db.backup-20261101-093312
created or upgraded database at /var/lib/ate/ate.db (schema version 5)
\end{lstlisting}

Three rules decide when a copy is made:

\begin{itemize}
  \item \textbf{A new database makes no copy} --- there is nothing to back up.
  \item \textbf{A database that needs nothing makes no copy}. A migration only runs when a table
    this ATE declares is missing, so a database already at the current schema is left alone and
    not copied.
  \item \textbf{A copy is never overwritten}. If that name is already taken, the next copy gets a
    counter: \texttt{\dots-2}, \texttt{\dots-3}.
\end{itemize}

The copy is taken with SQLite's backup API rather than \texttt{cp}, so it includes the writes that
are still sitting in \texttt{ate.db-wal}: a database with a Test Run in flight is copied whole and
consistent, which a file-by-file copy of \texttt{ate.db} alone would not be.

If the copy cannot be written --- the data directory is not writable, or another process is holding
the database --- the migration is refused and nothing is changed:

\begin{lstlisting}[language=ATEShell]
refusing to migrate: cannot write a backup of /var/lib/ate/ate.db to
/var/lib/ate/ate.db.backup-20261101-093312
nothing in the database was changed
\end{lstlisting}

A migration that starts and then fails leaves its copy where it is and prints the path; that copy is
how the upgrade is rolled back.

Copies are not removed automatically. They accumulate beside the database until the Administrator
deletes the ones they no longer need.

\section{What you still back up yourself}

The migration covers the database on the day it changes. It is not a backup policy, and it covers
nothing else:

\begin{itemize}
  \item \textbf{The config.} \texttt{config.json}. It is small, it is not backed up by anything, and
    \texttt{apt purge} deletes it on Linux.
  \item \textbf{The reports.} Test Reports are generated from the database, so a report can always be
    regenerated --- but only from a database that still exists. Treat the \texttt{reports} directory
    as the deliverables and the database as the archive; keep both.
  \item \textbf{The whole data directory, now and then.} The migration copy is only made when a
    migration runs. A database that never migrates never gets one.
\end{itemize}

The data directory is small enough to copy whole: \texttt{ate.db}, \texttt{reports}, \texttt{logs}
and \texttt{definitions}. Keeping \texttt{logs} is optional and is the only part of the data that
records ATE's own activity rather than a Test Run.

Restoring is the same copy, in the other direction. Stop ATE, put the copy at the path
\texttt{database.path} points to (with its \texttt{-wal} and \texttt{-shm} companions gone), and
start again: the schema check at launch accepts a database whose version matches this ATE, and the
runs it holds appear on the History screen as they did before. There is no partial restore and no
undo of a single run: a run is a row in \texttt{test\_run} with its steps in \texttt{node}, and the
only way back to an earlier state is the copy taken before the change.

\cando{check a database's version, migrate it safely, and say when ATE makes the copy, when it
does not, and what it still leaves to the Administrator.}
